\section{Magical Perspectives}

Updated to include the text, Evola's notes and a missing section. Note that the on-line texts available have several errors in them, which I have been tediously correcting. Please let me know if you find others.

\begin{quotationx}
In the final article of Book three of the Introduction to Magic, Julius Evola translates several sections from Aleister Crowley's Liber Aleph, the Book of Wisdom or Folly. Overlooking Crowley's grandiosity and deliberately cultivated ``bad boy'' image, which are understood as mere façades, Evola considers him to be of the first rank in the field of magic. I here translate Evola's introductory remarks, followed by the enumeration of the sections that he translated. The English text is online at: LIBER ALEPH VEL CXI\footnote{\url{http://www.rahoorkhuit.net/library/libers/pdf/lib\_0111.pdf}}.

Evola was still hung up on ``shocking the bourgeois''. However, by now in 2013, everything has lost its shock value. Sex and drugs are the preferred recreational activities, even, and perhaps especially, of the bourgeois. Nevertheless, the True Will is still elusive and the bourgeoisie are more confused than ever. A large portion of the population is barely functioning due to a variety of psychological and emotional problems\footnote{\url{http://www.webmd.com/mental-health/news/20040601/rate-of-mental-illness-is-staggering}}. Suicides are up\footnote{\url{http://www.nytimes.com/2013/05/03/health/suicide-rate-rises-sharply-in-us.html?\_r=0}}.

We can suppose the subtitle was deliberately chosen: it is a book of wisdom or folly, but who decides? It is surprising that Evola went along with the implicit egalitarianism of the maxim that every man is a star. That is both the common wisdom and the common folly of today. Of course, when women are now filling up facebook with quotes from Nietzsche, Crowley, etc., the impact is lost, but the consequences linger.

It was easy to shock the bourgeois by being outrageous at one time; there was actually little risk involved. Now there is danger in shocking anyone, although it can still be done. Perhaps the time has come for some new maxims, although they are not channeled from an Egyptian spook:

\begin{enumerate}
\item The world is hierarchically ordered
\item True Will lies in adhering to the cosmic law
\end{enumerate}

Nevertheless, Crowley was a knowledgeable Hermetist. I am interested in why Evola selected these particular passages, if anyone has any opinions.
\end{quotationx}

In the contemporary magical amphitheater, the Englishman Aleister Crowley is a figure of the first rank. According to the common judgment he belongs to the ``blacker'', explicitly satanic, direction, and he even desired this form of appearances, calling himself ``The Great Beast 666''. In reality, this was in a large measure a simple façade; by calling himself that, Crowley simply followed the ``Left Hand Path'', possessing a particular qualification for that. His life was extremely varied and complex: he was poet, painter, student of natural sciences, explorer, alpinist (also with two climbs in the Himalayas), in contact with secret organizations of every type; it is even said that he had had a secret part in some recent political movements.

He practiced ceremonial and evocative magic, and precisely in that way, he had transmitted in Cairo a type of revelation, or new law, from an entity called Aiwass, that he baptized the ``Law of Thelema'' and of which he made himself the promoter. Two maxims are its bedrock: ``Every man is a star'' and ``Do what thou wilt shall be the whole of the law''. Formulating the second maxim this way, Crowley, in line with his taste, wanted to `épater le bourgeois', because its effective meaning is not so much the following of unrestrained license, but rather to exclude every exterior and extraneous law and to assume one's deepest will, in respect to which every man appears to be the manifestation of a force from above (from a ``star''). Practically, Crowley's magic was close to Tantra: sexual magic and the use of drugs had a great part in it. It is said that different disciples of Crowley ended up insane or suicides (the two possibilities that always flank the ``direct path''); however, he himself died in 1947 at the age of seventy two, sane and in full possession of his faculties.

The passages that we have translated here belong to the Liber Aleph, the Book of Wisdom or Folly that at this time exists only as a manuscript. The reader, apart from a certain syncretism and a certain baroquism, will easily find here elements of traditional magical discipline.

Sections included are in this order:

\paragraph{On the training of the mind: Mathematics (47)}


Now concerning the first Foundation of thy Mind I will say somewhat. Thou shalt study with Diligence in the mathematics, because thereby shall be revealed unto thee the Laws of thine own Reason and the Limitations thereof. This Science manifesteth unto thee thy true Nature in respect of the Machinery whereby it worketh; and showeth in pure Nakedness, without Clothing of Personality or Desire, the Anatomy of thy conscious Self. Furthermore, by this thou mayst understand the Essence between the Relation of all Things, and the Nature of Necessity, and come to the Knowledge of Form. For this Mathematics is as it were the last Veil before the Image of Truth, so that there is no Way better than our Holy Qabalah, which analyseth all Things soever, and reduceth them to pure Number; and thus their Natures being no longer coloured and confused, they may be regulated and formulated in Simplicity by the Operation of Pure Reason, to thy great Comfort in the Work of our Transcendental Art, whereby the Many become One.

\paragraph{On the training of the mind: Classics (48)}


My son, neglect not in any wise the Study of the Writings of Antiquity, and that in the original Language. For by this thou shalt discover the History of the Structure of thy Mind, that is, its Nature regarded as the last term in a Sequence of Causes and Effects. For thy Mind hath been built up of these Elements, so that in these Books thou mayst bring into the Light thine own subconscious Memories. And thy Memory is as it were the Mortar in the House of thy Mind, without which is no Cohesion or Individuality possible, so that the Lack thereof is called Dementia. And these Books have lived long and become famous because they are the Fruits of ancient Trees whereof thou art directly the Heir, wherefrom (say I) they are more truly germane to thine own Nature than Books of Collateral Offshoots, though such were in themselves better and wiser. Yes, O my Son, in these Writings thou mayst study to come to the true Comprehension of thine own Nature, and that of the whole Universe, in the Dimension of Time, even as the Mathematic declareth it in that of Space: That is, of Extension. Moreover, by this Study shall the Child comprehend the Foundation of Manners: the which, as sayeth one of the Sons of Wisdom, maketh Man.

\paragraph{On the training of the mind: Science (49)}


Since Time and Space are the Conditions of Mind, these two Studies are fundamental. Yet there remaineth Causality, which is the Root of the Actions and Reactions of Nature. This also shalt thou seek ardently, that thou mayst comprehend the Variety of the Universe, its Harmony and its Beauty, with the Knowledge of that which compelleth it. Yet this is not equal to the former two in Power to reveal thee to thy Self; and its first Use is to instruct thee in the true Method of Advancement in Knowledge, which is fundamentally, the Observation of the Like and the Unlike. Also, it shall arouse in thee the Ecstasy of Wonder; and it shall bring thee to a proper Understanding of Art Magick. For our Magick is but one of the powers that lie within us undeveloped and unanalysed; and it is by the Method of Science that it must be made clear, and available to the Use of Man. Is not this a Gift beyond Price, the Fruit of a Tree not only of knowledge but of Life? For there is that in Man which is God, and there is that also which is Dust; and by our Magick we shall make these twain one Flesh, to the Obtaining of the Empery of the Universe.

\paragraph{On the virtue of daring (46)}


Yet this I charge thee with my Might: Live Dangerously. Was not this the Word of thine Uncle Friedrich Nietzsche? Thy meanest Foe is the Inertia of the Mind. Men do hate most those things which touch them closely, and they fear Light, and persecute the Torchbearers. Do thou therefore analyse most fully all those Ideas which Men avoid; for the Truth shall dissolve Fear. Rightly indeed Men say that the Unknown is terrible; but wrongly do they fear lest it become the Known. Moreover, do thou all Acts of which the common Sort beware, save where thou hast already full knowledge, that thou mayest learn Use and Control, not falling into Abuse and Slavery. For the Coward and the Foolhardy shall not live out their Days. Every Thing has its right Use; and thou art great as thou hast Use of Things. This is the Mystery of all Art Magick, and thine Hold upon the Universe. Yet if thou must err, being human, err by excess of courage rather than of Caution, for it is the Foundation of the Honour of Man that he dareth greatly. What sayth Quintus Horatius Flaccus in the third Ode of his First Book? Die thou standing!

\paragraph{On the key of dreams (17)}


And now concerning Meditation let me disclose unto thee more fully the Mystery of the Key of Dreams and Phantasies.

Learn first that as the Thought of the Mind standeth before the Soul and hindereth its Manifestation in consciousness, so also the gross physical Will is the Creator of the Dreams of common Men, and as in Meditation thou doest destroy every Thought by mating it with its Opposite, so must thou cleanse thyself by a full and perfect Satisfaction of that bodily will in the Way of Chastity and Holiness which has been revealed unto thee in thy Initiation.

This inner Silence of the Body being attained, it may be that the true Will may speak in True Dreams; for it is written that He giveth unto His Beloved in Sleep. Prepare thyself therefore in this Way, as a good Knight should do.

\paragraph{On the sleep of light (18)}


Now know this also that at the End of that secret Way lieth a Garden wherein is a Rest House prepared for thee. For to him whose physical Needs of whatever Kind are not truly satisfied cometh a Lunar or physical Sleep appointed to refresh and recreate by Cleansing and Repose; but on him that is bodily pure the Lord bestoweth a Solar or Lucid Sleep, wherein move Images of pure Light fashioned by the True Will. And this is called by the Qabalists the Sleep of Shiloam, and of this doeth also Porphyry make mention mention, and Cicero, with many other Wise Men of Old Time.

Compare, o my Son, with this Doctrine that which was taught thee in the Sanctuary of the Gnosis concerning the Death of the Righteous; and learn moreover that these are but particular Cases of an Universal Formula.

\paragraph{On the way of inertia (29)}


Of the Way of the Tao I have already written to thee, o my Son, but I further instruct thee in this Doctrine of doing everything by doing nothing. I will first have thee to understand that the Universe being as above said an Expression of Zero under the Figure of the Dyad, its Tendency is continually to release itself from that strain by the Marriage of Opposites whenever they are brought into Contact. Thus thy true Nature is a Will to Zero, or an Inertia, or Doing Nothing; and the Way of Doing Nothing is to oppose no Obstacle to the free Function of that true Nature. Consider the Electrical Charge of a Cloud, whose Will is to discharge itself in Earth, and so release the Strain of its Potential. Do this by free conduction, there is Silence and Darkness; oppose it, there is Heat and Light, and the Rending asunder of that which will not permit free Passage to the Current.

\paragraph{On the way of freedom (30)}


Do not think then that by Non-Action thou doest follow the Way of the Tao, for thy Nature is Action, and by hindering the Discharge of thy Potential thou doest perpetuate and aggravate the Stress. If thou ease not Nature, she will bring thee to Dis-Ease. Free thereof every Function of thy Body and of every other Part of thee according to its true Will. This also is most necessary, that thou discover that true Will in every Case, for thou art born into Dis-Ease; where are many false and perverted Wills, monstrous Growths, Parasites, Vermin are they, adherent to thee by Vice of Heredity, or of Environment or of evil Training. And of all these Things the subtlest and most terrible, Enemies without Pity, destructive to thy will, and a Menace and Tyranny even to thy elf, are the Ideals and Standards of the Slave-gods, false Religion, false Ethics, even false Science.

\paragraph{On the motion of life (20)}


Learn then, o my Son, that all Phenomena are the effect of Conflict, even as the Universe itself is a Nothing expressed as the Difference of two Equalities, or, an thou wilt, as the Divorce of Nuit and Hadit\textsuperscript{*}. So therefore every Marriage dissolveth a more material, and createth a less material Complex; and this is our Way of Love, rising ever from Ecstasy to Ecstasy. So then all high Violence, that is to say, all Consciousness, is the spiritual Orgasm of a Passion between two lower and grosser Opposites. Thus Light and Heat result from the Marriage of Hydrogen and Oxygen; Love from that of Man and Woman, Dhyana or Ecstasy from that of the Ego and the non-Ego.

But be thou well grounded in this Thesis corollary, that one or two such Marriages do but destroy for a Time the Exacerbation of any Complex; to deracinate such is a Work of long Habit and deep Search in Darkness for the Germ thereof. But this once accomplished, that particular Complex is destroyed, or sublimated for ever.\footnote{Evola's note: Nuit and Hadit are, in the terminology used by C., with a certain reference to the ancient Egyptian tradition, the cosmic Feminine and Masculine, the ``nothing'' or Zero, the equivalent of ``emptiness'', \emph{sunya}, of Mahayana Buddhism.}

\paragraph{On Diseases of the Blood (21)}


In the text, this section was joined to the preceding. The emphasis at the end was added by Evola.

Now then understand that all Opposition to the Way of Nature createth Violence. If thine excretory System do its Function not at its fullest, there come Poisons in the Blood, and the Consciousness is modified by the conflicts or Marriages between the elements heterogeneous. Thus if the Liver be not efficient, we have Melancholy; if the Kidneys, Coma; if the Testes or Ovaries, loss of Personality itself. Also, an we poison the Blood directly with Belladonna, we have Delirium vehement and furious; with Hashish, Visions phantastic and enormous; with Anhiolonium, Ecstasy of colour and what not; with diverse Germs of Disease, Disturbances of Consciousness varying with the Nature of the Germ. Also with Ether, we gain the Power of analysing the Consciousness into its Planes; and so for many others.

But all these are, in our mystical Sense, Poisons; that is, we take two Things diverse and opposite, binding them together so that they are compelled to unite; and \emph{the Orgasm of each Marriage is an Ecstasy, the Lower dissolving in the Higher}.

\paragraph{On wisdom in sexual affairs (44)}


Consider Love. Here is a Force destructive and corrupting where by many Men have been lost. Yet without Love Man were not Man. Therefore thine Uncle Richard Wagner made of our Doctrine a musical Fable, wherein we see Amfortas, who yielded himself to Seduction, wounded beyond Healing; Klingsor, who withdraw himself from a like Danger, cast out for ever from the Mountain of Salvation; and Parsifal, who yielded not, able to exercise the true Power of Live, and thereby to perform the Miracle of Redemption. Of this also have I myself written in my Poema called Adonis. It is the same with Food and Drink, with Exercise, with Learning itself, the Problem is ever to bring the Appetite into right Relation with the Will. Thus thou mayst fast or feast; there is no Rule than that of Balance. And this Doctrine is of general Acceptation among the better Sort of Men; therefore on thee will I rather impress more carefully the other Part of my Wisdom, namely, the Necessity of extending constantly thy Nature to new Mates upon every Plane of Being, so that thou mayst become the perfect Microcosm, an Image without Flaw of all that is.

\paragraph{On the mystical marriage (23)}


O my Son, how wonderful is the Wisdom of this Law of Love! How vast are the Oceans of uncharted Joy that lie before the Keel of thy Ship! Yet know this, that every Opposition is in its Nature named Sorrow, and the Joy lieth in the Destruction of the Dyad. Therefore, must thou seek ever those Things which are to thee poisonous, and that in the highest Degree, and make them thine by Love. That which repels, that which disgusts, must thou assimilate in this Way of Wholeness. Yet rest not in the Joy of the Destruction of each complex in thy Nature, but press on to that ultimate Marriage with the Universe whose Consummation shall destroy thee utterly, leaving only that Nothingness which was before the Beginning.

So then the Life of Non-Action is not for thee; the Withdrawal from Activity is not the Way of the Tao; but rather the Intensification and making universal every Unit of thine Energy on every Plane.

\paragraph{A final thesis on love (6)}


Therefore, o my Son, be thou wary, not bowing before the false Idols and ideals, yet not flaming forth in Fury against them, unless that be thy Will.

But in this Matter be prudent and be silent, discerning subtly and with acumen the nature of the Will within thee; so that thou mistake not Fear for Chastity, or Anger for Courage. And since the fetters are old and heavy, and thy Limbs withered and distorted by reason of their Compulsion, do thou, having broken them, walk gently for a little while, until the ancient Elasticity return, so that thou mayst walk, run, and leap naturally and with Rejoicing.

Also, since these Fetters are as a Bond almost universal, be instant to declare the Law of Liberty, and the full Knowledge of all Truth that appertaineth to this Matter; for if in this only thou overcome, then shall all Earth be free, taking its Pleasure in Sunlight without Fear or Phrenzy. Amen.

\paragraph{On the bull (153)}


Concerning the Bull, this is thy Will, constant and unwearied, whose Letter is Vau, which is Six, the Number of the Sun. He is therefore the Force and the Substance of thy Being; but besides this, he is the Hierophant in the Taro, as if this were said: that thy Will leadeth thee unto the Shrine of Light. And in the Rites of Mithras the Bull is slain, and his Blood poured upon the Initiate, to endow him with that Will and that Power of Work. Also in the land of Hind is the Bull sacred to Shiva, that is God among that Folk, and is unto them the Destroyer of all Things that be opposed to Him. And this God is also the Phallus, for this Will operateth through Love even as it is written in our Own Law. Yet again, Apis the Bull of Khem hath Kephra the Beetle upon His tongue, which signifieth that it is by this Will, and by this Work, that the Sun cometh unto Dawn from Midnight. All these Symbols are most similar in their Nature, save as the Slaves of the Slave- gods have read their own Formula into the Simplicity of Truth. For there is naught so plain that Ignorance and Malice may not confuse and misinterpret it, even as the Bat is dazzled and bewildered by the Light of the Sun. See then that thou understand this Bull in Terms of the Law of this our Æon of Life.

\paragraph{On the lion (154)}


Of this, Lion, o my Son, be it said that this is the Courage of thy Manhood, leaping upon all Things, and seizing them for their Prey. His letter is Teth, whose Implication is a Serpent, and the Number thereof Nine, whereof is Aub, the secret Fire of Obeah. Also Nine is of Jesod, uniting Change with Stability. But in The Book of Thoth He is the Atu called Strength, whose Number is E L E V E N which is Aud, the Light Odic of Magick. And therein is figured the Lion, even T H E B E A S T, and Our Lady B A B A L O N astride of Him, that with her Thighs She may strangle Him. Here I would have thee to mark well how these our Symbols are cognate, and flow forth the one into the other, because each Soul partaketh in proper Measure of the Mystery of Holiness, and is kin with his Fellow. But now let me show how this Lion of Courage is more especially the Light in thee, as Leo is the House of the Sun that is the Father of Light. And it is thus: that thy Light, conscious of itself, is the Source and Instigator of thy Will, enforcing it to spring forth and conquer. Therefore also is his Nature strong with hardihood and Lust of Battle, else shouldst thou fear that which is unlike thee, and avoid it, so that thy Separateness should increase upon thee. For this Cause he that is defective in Courage becometh a Black Brother, and to Dare is the Crown of all thy Virtue, the Root of the Tree of Magick.

\paragraph{Further on the lion (155)}


Lo! In the firs of thine Initiations, when first the Hoodwink was uplifted from before thine Eyes, thou wast brought unto the Throne of Horus, the Lord of the Lion, and by Him enheartened against Fear. Moreover, in Minutum Mundum, the Map of the Universe, it is the Path of the Lion that bindeth the two Highest Faculties of thy Mind. Again, it is Mau, the Sun at Brightness of high Noon, that is called the Lion, very lordly, in our Holy Invocation. Sekhet our Lady is figured as a lioness, for that She is that Lust of Nuit toward Hadit which is the Fierceness of the Night of the Stars, and their Necessity; whence also is She true Symbol of thine own Hunger of Attainment, the Passion of thy Light to dare all for its Fulfilling. It is then the Possession of this Quality which determineth thy Manhood; for without it thou art not impelled to Magick, and thy Will is but the Salve's Endurance and Patience under the Lash. For this Cause, the Bull being of Osiris, was it necessary for the Masters of the Æons to incarnate me as (more especially) a Lion, and my Word is first of all a Word of Enlightenment and of Emancipation of the Will, giving to every Man a Sprint within Himself to determine His Will, that he may do that Will, and no more another's. Arise therefore, o my son, arm thyself, haste to the Battle!

\paragraph{On the man (156)}


Learn now that this Lion is a natural Quality in Man, and secret, so that he is not ware thereof, except he be Adept. Therefore is it necessary for thee also to know, by the Head of the Sphinx. This then is thy Liberty, that the Impulse of the Lion should become conscious by means of the Man; for without this thou art but an Automaton. This Man moreover maketh thee to understand and to adjust thyself with Environment, else being devoid of Judgment, thou goest blindly upon an headlong Path. For every Star in his Orbit holdeth not his Way obstinately, but is sensitive to every other Star, and his true Nature is to do this. Oh how many are they whom I have seen persisting in a fatal Course, in Sway of the Belief that their dead Rigidity was Exercise of Will. And the Letter of the Man is Tzaddi, whose Number is Ninety; which is Maim, the Water that conformeth itself perfectly with its Vessel, that seeketh constantly its Level, that penetrateth and dissolveth Earth, that resisteth Pressure maugre its Adaptability, that being heated is the Force to drive great Engines, and being frozen breaketh the Mountains in Pieces. O my Son, seek well To Know!


\flrightit{Posted on 2013-08-20 by Cologero}

\begin{center}* * *\end{center}

\begin{footnotesize}
\begin{sffamily}
\texttt{cosmodromium on 2013-08-21 at 00:03 said:}

\url{http://hermetic.com/crowley/book-of-wisdom-or-folly/}


\hfill

\texttt{Bob on 2013-08-21 at 00:13 said:}

``every man or woman is a star''... is this really egalitarian?

Not all stars burn in equal brightness. I think the main point is that everyone has a true course ( a True Will) in their life that fits their nature, just as each star has its proper place and course in the heavens.

Its not that everyone is equal, but that everyone has their proper place.

\hfill

\texttt{B.M.M. Reynolds on 2013-08-21 at 08:38 said:}

Is the Law of Thelema compatible with the revealed truth of Christianity?

\hfill

\texttt{Deetaleh on 2013-08-21 at 09:59 said:}

Liber Aleph is certainly not one of the ``easiest'' or the most accessible works of Crowley. I think Evola simply did choose those passages because of their simplicity and practical value for the practitioner. Liber Aleph has quite a lot of Thelemic symbolism and concepts that lot of readers probably would have hard time comprehending.

I also do not think that even Evola had read Crowley/Thelema/Book of the Law, enough to make the most sense out of every passage.

\hfill

\texttt{JA on 2013-08-21 at 11:49 said:}

I have a feeling that Evola, a man who was uninitiated himself despite his best efforts to try to do it himself plus looking for any group he could find to the point of trying to get information out of Guenon, had too little powers of discernment when it came to people who actually had spiritual contacts with the higher realms, Evola tended to endorse anyone who could prove they actually had achieved contact with spiritual reality and he did not look deeper. Thus Evola endorsed Gurdjieff, whereas Guenon had feared him. Both were wrong as Gurdjieff did have some knowledge that he distorted and perverted, as Mouravieff shows -- but the point is that Evola grasped for any straw he could find. I believe that is the case with Crowley. Not to also mention that Crowley's personality was very very very akin to Evola's.

Of course the modern exponents of Crowley include mostly rock musicians, drug addicts, sexual perverts, and transsexuals (personal experience with all of the above).

\hfill

\texttt{Deetaleh on 2013-08-21 at 12:41 said:}

JA wrote: ``Evola tended to endorse anyone who could prove they actually had achieved contact with spiritual reality and he did not look deeper. Thus Evola endorsed Gurdjieff, whereas Guenon had feared him. Both were wrong as Gurdjieff did have some knowledge that he distorted and perverted, as Mouravieff shows -- but the point is that Evola grasped for any straw he could find.''

Crowley also had met Gurdjeff in person, but didn't think that highly of him even though he writes in his diaries: ``Gurdjieff, their prophet, seems a tip-top man. Heard more sense and insight than I've done for years. Gurdjieff clearly a very advanced adept''. There is also some rumors of some minor `verbal conflict' between the two on their parting, but they are probably highly exaggerated stories,

Crowley certainly was a complex figure and even after years of studying his work I really cannot make up my mind on him, but his various claims for certain `attainments' should not be taken as some idle talk. He was no charlatan.

\hfill

\texttt{h.ontologia on 2013-08-21 at 18:54 said:}

Very very very fearful you seem JA.

Let us consider Matthew 7:2.

\hfill

\texttt{Pickman on 2013-08-22 at 08:43 said:}

Would it be possible for Gornahoor to run a series on ``Crowley viewed from the Right'' similar of course to the Evolian review?

I understand that it may be a lot to ask for considering the breath of the undertaking but it would be much appreciated by the readers here and it would certainly help for future reference points to those lost in the modern melee.

\hfill

\texttt{JA on 2013-08-22 at 10:14 said:}

I am not fearful, rather the opposite. I used to believe in and live Thelema.

@ Mr Reynolds, Tomberg can help reconcile thelema with christ esoterically. But outwardly no if you believe Crowley was a prophet and The Book of the Law is inspired you have to accept the anti-Christian polemic contained in that book.

Crowley was not a fake, indeed, my question is -- which side was he on ? I do not know.

Pretty much everything Crowley offers can be found in Mouravieff- including sex magic (the fifth way).

\hfill

\texttt{Deetaleh on 2013-08-22 at 10:33 said:}

For those who may be interested: Crowley's magical motto for the grade of Adeptus Minor (5=6) was `Christeos Lucifitias', meaning more or less ``Christ-like Brightness''

It is interesting to note that Crowley never did himself talk about this motto and kept quiet about it. Even Regardie did not know his motto for the grade and his best guess was ``Heart of Jesus Girt About By a Serpent''

It was later revealed to be `Christeos Lucifitias' by Kaczynski who had full access to Crowley's material.

\hfill

\texttt{Acratophorus on 2013-08-22 at 18:27 said:}

Thanks for this. It was a fun read.

There are already a few of us attempting to look at Crowley / Thelema from a traditionalist perspective. We believe his work shows genuine contact with transcendent forces and, despite his personal failings, his legacy shows the most promise of providing a foundation for the real re-emergence of the perennial tradition in the West. This work is both intellectual and experimental, and has been going on now for some time among a small group of dedicated individuals.

We have not been very public about it as, over the last decade, we were making a systematic study of the existing Thelemic and pseudo-Thelemic organizations (Wicca, Satanism, post-apocalyptic Paganism, etc..) in order to determine the extent to which they were still salvageable, or if they have entirely fallen prey to the process of ``occult infiltration''. Sadly, the later seems to be the case, although some of us still feel that fighting a ``holding action'' against the nihilism of the Kali Yuga is valuable as a personal gesture of defiance.

Many of us enjoy your blog a great deal, and look forward to your continued work.

\hfill

\texttt{Cologero on 2013-08-22 at 20:14 said:}

Interesting point, Bob. There are black stars, dwarf stars, red giants, supernovas that burn brightly for a moment and then fade, bright stars with a steady light, and even a few starts to guide your navigation. What sounds like an egalitarian complement is really quite mischievous!

Mr. Reynolds, I don't want to try to make that case. On the other hand, St Augustine said, ``Love and do what thou wilt,'' assuming your ``True'' will. Guenon claims that Gargantua and Pantagruel is an esoteric text; Crowley adopted many ideas from it.

Deetaleh, so we still don't know why those passages were chosen, and in that order. Perhaps Evola was just showing off or perhaps they need a deeper look.

JA, that was the point of the post. Crowley could not be considered outrageous today, since his morals are the same as a PTA mom in the suburbs, never mind decadent rock stars. OTOH, Hermetists are often tricksters (perhaps I will explain why some day); e.g., look up Cagliostro on Gornahoor.

Pickman, I have considered that, but the only suitable work for such a project would be the Book of the Law. However, there is an admonition in it to not produce such a commentary; I take that seriously.

Deetaleh, in the traditional Catholic calendar, one of the days has a prayer to Lucifer.

Acratophorus, like you, I have applied traditional metaphysics to Tomberg's work; it grounds things in a better way and makes things more comprehensible. The task is to enrich the Hermetic Tradition. Tomberg advises, in this regard, to act in the spirit of free enquiry and to respect the tradition. That is why I provide difficult to find material, without enforcing particular interpretations.

Nevertheless, there is a choice to be made, viz., between the English path or the French path. When certain ideas enter the public domain, as has happened with the Order of the Golden Dawn and its derivatives like the OTO, distortion results and the results are harmful. Anyone interested in the topic of the future course of the Hermetic Tradition should read my commentary and retranslation of \textit{Tomberg's Preface to MOTT}\footnote{\url{http://www.meditationsonthetarot.com/enriching-the-tradition}}. I'm interested in hearing the pros and cons.

\hfill

\texttt{francismercuri on 2013-08-24 at 23:30 said:}

Guenon wrote that every being is necessary in that they ``hold their place'' in the cosmos: every being that exists (or existed), exists because they are a possibility in infinite possibility, were they not a possibility--which means a unique possibility in infinite possibility--they would not nor ever could exist, which is obvious. That they exist or existed evidences that they manifested some unique possibility, which because it exists, or happens to be a possibility, is now a necessity, because to deny it, or ``remove '' it from existence (really Being), once it has specifically been, would enter a limit upon Universal Possibility, a contradiction--so, every being ``holds their place''. The dictum ``Every man and woman is a star'' could very well be interpreted similarly--the brightness and lives of stars has been mentioned above, to which we can add the symbolism of constellations. While having individual characteristics, the stars too ``hold their place'' in ``families''--constellations--which in the traditional sciences of course, are characterized by distinct qualities, natures, and rulerships--with each individual star (possibility) contributing to a synergy. A more suitable ``egalitarian'' cliche would have been ``Every man and woman is a Planet (``wanderer''), not a star--which is astrologically speaking ``fixed'' in the celestial sphere, and moving as a group.

\hfill

\texttt{anon on 2013-08-25 at 07:41 said:}

When thinking about the intersection of Trad, magic and Christianity, I tend to think of Lisiewski. I read a little bit of his evocation book, and have seen Guenon mentioned when Lisiewski is discussed. Don't know enough about Joseph to form a solid opinion, but he's an interesting critter.

\url{http://www.occultforum.org/forum/viewtopic.php?f=8\&t=20197}


\hfill

\texttt{JA on 2013-08-26 at 10:50 said:}

I have never read Lisiewski before but just googled him... .and am going to read the whole book... ... ..as a Catholic he may be offering just what I am looking for... ... ... 

\hfill

\texttt{anon on 2013-08-26 at 11:01 said:}

JA, you might like this one:

\url{http://www.scribd.com/doc/154996540/Howlings-From-the-Pit}


I was skimming it yesterday, and I liked it a lot better than his other book that I had skimmed. The part about dumping your New Age books and paraphernalia, in order to free up trapped energy, which I've seen him write about before, repeatedly, is simply awesome. And applies to much more than just dumping New Age books and paraphernalia. Getting rid of ``spiritual clutter''.

\hfill

\texttt{JA on 2013-08-26 at 13:03 said:}

There is little information to be found on Lisiewski, he mysteriously retired in 2008 after writing his 3 books... ... .this is very interesting... ... people debate whether or not Joseph Lisiewski was even his real name... ... .I wonder if he MAY have been possibly a genuine initiate/magician... ... .

\hfill

\texttt{anon on 2013-08-26 at 13:47 said:}

Some people seem to think that he's Hyatt, since Hyatt died in 2008. But the style and content seem too different to me. More than 3 books though, mon ami. At least two novels on top of that, and a book about Regardie. No idea if he is for real. Some of his content is gold.

\hfill

\texttt{Cologero on 2013-08-27 at 22:44 said:}

Acratophorus, if you engage in a Hermetic conversation, you start on common ground and move deeper. Nothing I wrote is ``simply incorrect'' and is not based on mindless suppositions, ``grave errors'', or ``profound misstatement of the facts''. I'll attribute your ignorance to your lack of familiarity with PTA ladies and with what they are willing to do when properly coaxed. I recommend a different approach: instead of attacking, seek first to understand the point and not react mechanically to what you think I meant.

I had in mind C's defense of sexual expression, particularly homosexuality, in the Book of the Law. There is little that is objectionable to the educated classes today, and to hold the contrary position is considered beyond the pale. Is that really ``simply incorrect''? Perhaps there may be nuances, but the essential point is clear. To move forward, as C himself makes clear, whatever he writes needs to be understood on the level it is intended, whether the general public or a particular degree of attainment. In particular, coprophagia was never intended for the public, nor part of his social vision, if it could be said he had one, but was a test for a certain level. If all is samsara, then one should be indifferent to the particular experience, but this is not the place to expand on that. As a matter of fact, we have quoted the \textit{Dalia Lama who said}\footnote{\url{http://www.gornahoor.net/?p=3356}} that a meal of crap with a glass of pus should be as appetizing to us as the finest meal. This level of indifference is a prerequisite to Buddhist Tantric Initiation. So your assumption that there is something new under the sun, and that you have discovered it, is simply incorrect.

I didn't realize I called Crowley a ``poseur'' as you seem to believe, so I did a text search and did not find the word anywhere. Unless you can locate the context of that claim, I expect you to retract that accusation. It is true that Evola referred to Crowley's ``facade'' that hid the deeper aspects of his thought by repulsing certain people. I agree with Evola in that regard, so I am not really interested in that outer facade; nevertheless, Crowley did make significant contributions to Hermetism. I suspect, on the other hand, that many are attracted to Crowley precisely because of the facade, and not for much else.

Since my skin infection is subsiding, I am able to comment a bit on Liber Aleph rather than continue to wait for George to do it. Maybe there will be a few who will read it in the right spirit.

\hfill

\texttt{francismercuri on 2013-08-29 at 15:09 said:}

What is a ``real magician'' and a ``real initiate'', JA? I'd not necessarily conflate the two. What some might find as a useful to meditate is the idea that the Western magical/grimoiric tradition, as exemplified and collected by Agrippa for one, is aimed not at ``astral magic'' but Gnosis, through techniques and ``purifications''--exercises--mostly plainly laid out in the texts:\\
\url{http://www.academia.edu/1170523/Better\_than\_Magic\_Cornelius\_Agrippa\_and\_Lazzarellian\_Hermetism\_2009\_}


If Thelema and Christianity are ``compatible'' is open to endless debate (which I don't intend to engage). However, Crowley had written that he desired to restore Christianity to its ``Solar-Phallic'' character (although in terms of epoch and surrounding traditions, one can understand the classification of ``Solar-Phallic''--Christianity imo, recapitulates the Polar symbolism as well--something Crowley either ignored or could not see)--a comment that has lead some to regard Thelema as an ``Christian esotericism''. Nonetheless, regardless of one's feelings about Thelema, it ought not be overlooked that his A.:A.: curriculum and goal driven method of advancement is one of the best organized by Orders openly publishing their materials, and is more or less consistent with many of the ideas contained in the above Haangraaf paper.

Rabelais' name has also been associated with a short tract on Spirit evocation:\\
\url{http://www.sacred-texts.com/grim/moses7/m7112.htm}


\hfill

\texttt{francismercuri on 2013-08-29 at 17:20 said:}

More regarding the A.:A.: curriculum can be found and examined in ``Liber Zelotes'' which includes a listing of prospective A.:A.: Grade examinations. It seems important to look upon these aims and skills, not as ``Thelemic'', but what they are in themselves, why one would do them, and what they might lead to in themselves. In fact, many of the nuts and bolts core techniques aren't too different than some discussed on Gornahoor:

\url{http://www.scribd.com/doc/40755160/Liber-Zelotes}


This might be compared/contrasted or worked along with something like the ``tasks'' assigned in contemporary Platonic ``Orders'':

\url{http://www.openingmind.com/aboutus.html}


\hfill

\texttt{francismercuri on 2013-08-29 at 23:38 said:}

Cologero, an anecdote: Isreal Regardie in explaining his reasoning behind the selections from Crowley's ``Equinox'' selected to compose ``Gems From The Equinox'', explains excluding Crowley's translation of Eliphas Levi's ``Key of the Mysteries'' because Regardie ``cannot find anything useful'' in the ``French school of occultism''.

``Pros and cons'' regrading a ``choice'' of English or French Hermeticism. In ``Enriching the Tradition'' it seems that the parameters defining ``English Hermeticism'' are Golden Dawn and offshoots. I would ask though why begin with Golden Dawn? For instance, ``Theatrum Chemicum Brittanicum'' can be regarded as a sort of original English Hermetic scripture, and a very high level of Hermetic expression at that, in Bardic verse. This branch of tradition is quite compact, but very thick, and includes more so than the Golden Dawn side, Dr. Dee's works. Thomas Vaughn who also precedes the GD (although his works were somewhat employed by GD, as in Westcott's inclusion of ``Euphrates'' in ``Collectanea Hermetica''), represents another branch of this English Hermeticism worth reviving.

In terms of Sacred Sciences the English tradition has also included astrologers--whose writings express traditional cosmology, and directly reference a ``Hermetic'' origin/source, such as William Lilly, William Ramesy; astro-physicians William Saunders, Nicholas Culpeper, and even recently Dr. H.L. Cornell who certainly greatly preserved so very much of traditional Western/English medicine in ``Encyclopedia of Medical Astrology'' (which, interestingly, in the 1930's when Cornell, an MD wrote/practiced, this astrology was still taught in medical schools--``establishment'' schools,not indie colleges). Although debate exists surrounding who gets credit for ``inventing'' calculus, the role of Newton as introducing to the English at the very least is obvious--and this, as Guenon wrote, is an extremely valuable metaphysical language, as all number.

Another important ``current'' in pre-GD ``English'' (now Celtic) tradition would of course include Eriugena, who links the Celtic tradition with that of Dionysius the Aeropagite and Platonism. Can we perhaps see the transmission of certain ideas too in William Law with Boehme, and Thomas Taylor with Platonism? It would seem that even with Pound there is more than meets the eye:

\url{http://koha.arcadia.edu/cgi-bin/koha/opac-detail.pl?biblionumber=18123&shelfbrowse\_itemnumber=20043}

Additionally, with Shakespeare, over recent years a flurry of esoteric interpretations have emerged, connecting his writings with all sorts of ideas--Martin Ling's of course presents the most ``traditional'' esotericism in this area.

I'm sure others could add to the list. So, this was neither pro nor con. Rather, I propose to simultaneously develop this list with other inclusions/exclusions, to help define parameters describing English Hermeticism--which would better help in considering whatever pros or cons are before us. For example, do any of these inclusions stray to far afield of what we wish to describe as ``Hermeticism''?

\hfill

\texttt{francismercuri on 2013-08-30 at 00:16 said:}

An important ``mandalla'' of English Hermeticism:

\url{http://www.levity.com/alchemy/rscroll.html}

\url{http://www.levity.com/alchemy/ripscrl1.html}


In my opinion, the Ripley Scroll alone offers an incredible amount to meditate and ``unpack'' when considering English Hermeticism. This scroll is presented as a single emblem of the Great Work by Ripley, with all ingredients and through each stage of the process--it can also be seen as emblem of all his verse. It might perhaps be seen as English equivalent of the famous Taoist Inner Alchemy scroll, ``The Neijing Tu'':

\url{http://en.wikipedia.org/wiki/File:Neijingtu\_from\_KamLanKoon.jpg}


This article is a general survey of pre-GD English Hermeticism:

\url{http://en.wikipedia.org/wiki/File:Neijingtu\_from\_KamLanKoon.jpg}

Tillyard's book ``Elizabethan World Picture'' is scholarly source describing the presence of a mostly Medieval world view in Elizabethan England, and where the most pronounced expression of those spiritual/esoteric aspects can be located. As an aside, this book is on the ``required reading'' list of certain practical astrological courses:

\url{http://www.amazon.com/Elizabethan-World-Picture-Vintage/dp/0394701623/ref=sr\_1\_1?s=books\&ie=UTF8\&qid=1377835806\&sr=1-1\&keywords=elizabethan+world+picture+by+tillyard+9780394701622}

\hfill

\texttt{francismercuri on 2013-08-30 at 05:11 said:}

Cologero: I hope what I'm doing does not seem too far afield of Evola and Liber Aleph, as it has a relevance to Crowley and this thread in general, but also a connection with the ``Enriching the Tradition'' post.

The dichotomy presented was more or less contemporary English vs. French Hermeticism, with the English represented by GD, Crowley, and OTO, and treated as a pagan branch; while on the French side Tomberg is considered a recent chief representative of a very lengthy and colorful branch of Christian Hermeticism. I've previously mentioned the pre-GD English Hermetic/Alchemical (or Platonic) components (particularly the ``Theatrum'' and Ripley emblems) that I consider anterior to, and ``purer'' than the GD and related ``currents''--which is overwhelmingly Christian. Now, another significant contemporary representative of Hermeticism that ought not be overlooked is part of neither the English nor French schools, but the German/Czech. Franz Bardon adds another variable in our overall Hermetic equation, and while contemporaneous with the GD era, is like Tomberg, Christian (Bardon often speaks in Christian theological language, and in keeping reverence, of ``Our Lord Jesus'').

Years ago, a very conservative ``Guenonist'' repudiated and dismissed Bardon as yet another ``occultist'' to be shunned. Or, at least not taken seriously. Yet, Bardon deserves to be taken seriously as major contemporary representative of German/Czech Hermeticism:

1. Bardon's three major works are considered as developments of the first three trumps of the Tarot. ``The Magician'' corresponds with ``Initiation Into Hermetics'', ``The High Priestess'' corresponds with ``Practice of Magical Evocation'', and ``The Empress'' with ``Key to the True Qabbalah''.

2. Bardon's Hermeticism, even if he was not associated with a ``regular'' Order, had a living traditional component in the background via his training in medicine. Bardon graduated from a naturopathic/homeopathic medical college--he was also comprehensive in at least medical spagyrics, as he produced his own ellirs according to those methods. This is telling, because at that time, such schools, especially in his nation, would have been studying Hahnemann directly in ``Organon'', as well as the Paracelsian canon. It must be remembered that contemporary spagyrics owes much to Paracelsus, as does Hahnemann, whose work/school is not ``Paraclsian'' per se, but a development of certain tenets (actually, another stone that has been insufficiently turned involves fully reading ``Organon'' as metaphysical text--this due mostly to its style of aphorism and maxim, which are describing natural laws). If anything nothing else then ``exteriorly'', Bardon received Hermetic Initiation via his Paracelsian/Hahnemannian/Spagyric instructors.

3. Despite how much Bardon's work seems directed toward and obsessed with ``powers'', a ``vichara'' and techniques to attain vichara fuel most of his books basic premises concerning man.

4. A similarity can be seen between Bardon's treatment of the ``weapons'' of magic as symbols of internal states and powers, with Haangraff's treatment of Agrippa, and Crowley's treatment of the same implements and accoutraments in ``Book 4''.

5. Bardon did discernibly have access to primary Grimoiric sources, and an understanding of them--which is mostly what he is doing--restating the Grimoires as internal exercises/experiements. As an example though, of reliance on primary sources, in ``Evocation'' the names of the Lunar beings are in fact anagrams of the names in ``The Picatrix''--Bardon lists them all according to Lunar Mansion degree and minute--and the correspondences line up; clearly, his system was not an modern ``invented'' occultism, but contained quite a bit coming straight out of several appliacations of Hermeticism.

6. This is just a curiosity: In early editions of ``Inititation Into Hermetics'' (not sure if new eds have it), early in the text, Bardon completes a comment with ``Love under Will'', in quotes. Was this connecting him with Crowley or Rabelais, or both? Was this a nod toward Crowley despite Bardon's Christianity--had Bardon even read Crowley--this I don't know.

7. Michael Maier is a Gnostic Saint in Crowley's Gnostic Church--more on that later... .

\hfill

\texttt{francismercuri on 2013-08-30 at 06:37 said:}

Before I forget, let me include an addendum to the ``German/Czech Hermeticism'' grouping. I've not finished this anthology, but some of the considerations, nuance our understanding of an Polish, Slovakian, Hungarian, Moravian, Eastern European Hermeticism and magic:

``The Role of Magic in the Past''

\url{http://www.forumhistoriae.sk/e\_kniznica/szeghyova.pdf}


Aside from the more obvious picks such as Boehme and Eckhardt, and the ``contemporary'' Hermeticism of a Bardon, or Christian esotericism of Steiner, this ``school'' would also Hermetically include the schools of men like Basil Valentine, Sendivogious, and Maier, as alchemical adepts, and others like Johannes Trithemius as the Grimoiric/Theurgic input.

\hfill

\texttt{Cologero on 2013-08-30 at 08:08 said:}

francis: This French Martinis site, \textit{Ordre Martiniste Opératif Au Québec}\footnote{\url{http://www.martinismeoperatifquebec.org/indexeng.html}}, accepts Franz Bardon, but I don't know if he ever had any relationship to any such order. I did not find that quote in the text you mentioned.

Tomberg mentioned a ``living'' tradition since the 18th century, so I don't know if the English figures you mention qualify. Tomberg doesn't seem to reference them, although Steiner did in some cases. It may be that Tomberg wanted to focus on one particular stream, the one that he knew best, and perhaps it does not represent a rejection. Tomberg does indeed mention Crowley and Westcott in passing a few times, so he was not unfamiliar with that movement. Nevertheless, to ``enrich the tradition'' certainly does not mean to keep it static, and restricted to certain figures.

\hfill

\texttt{francis mercuri on 2013-08-30 at 16:04 said:}

In ``Initiation Into Hermetics'', ``Step X''-``Magical Mental Training'', Bardon is discussing interactions with water spirits, and the magical danger of becoming infatuated with them, advising the student that ``He must keep the motto in mind: love is the law, but love under a strong will''.

\hfill

\texttt{francismercuri on 2013-08-31 at 23:47 said:}

More concerning Bardon's quote. Not sure what this could lead to (if anything)--but, it is curious that Bardon works the ``Love is the Law'' quote into his section concerning water spirits, mermaids. The form the quote takes in Bardon does not sound like St. Augustine's rendition. This is quite interesting as these mermaids relate to the Melusine legend(s)--one of which has Melusine involved in constructing the historical Abbey of Maillezais--of which Rabelais was inspired, and it is in Maillezais, that Rabelais becomes a Benedictine.

\hfill

\texttt{anon on 2013-09-08 at 20:16 said:}

\textit{Robert Anton Wilson on The Holy Guardian Angel}\footnote{\url{https://www.youtube.com/watch?v=n7QhT39NBuE}}

NOTE from moderator. Anon, we don't particularly value unattributed links to other sites. There should be some ``added value'', e.g., what in particular is of note in the link, what point does it illustrate, etc. Henceforth, all such comments will be rejected.

I can point out to those tempted to follow the link that is a truly boring and purposeless talk and has absolutely nothing to do with the alleged topic.

\hfill

\texttt{anon on 2013-09-10 at 13:11 said:}

Some things about Wilson I truly love. In others spheres, I find him intellectually cowardly. The video? I remember him once mentioning a laundry list of things that he felt would be useful to someone getting involved with magic. Among these, he mentioned Alfred Korzybski's general semantics. As your more discerning readers (who watched the vid!) have no doubt already picked up on, the video is general semantics applied to the K\&C of the HGA, specifically. And to psi/occult phenomena in general.

We are playing (on this blog) with terms and concepts. Playing at being pundits. Some here might be actual baseball players, while others may be more like baseball card collectors. Memorizers and reciters of statistics. Dropper of names. Masturbators of the ego.

A young black guy from inner-city Philly, over a decade ago, first introduced me to the ball player/card collector meme. And I cannot shake it, so helpful do I find it.

Fighting and arguing on the Internet about baseball statistics is okay, but it is not for everyone. Some people also like to play baseball.

\hfill

\texttt{Cologero on 2013-09-10 at 16:17 said:}

Yes, anon, and what do we do about the people who don't even realize there is a baseball game going on?

\hfill

\texttt{anon on 2013-09-10 at 16:23 said:}

Life is too short.

\hfill

\texttt{Cologero on 2013-09-11 at 00:17 said:}

Life is short, anon, and art is long; hence the show must go on. A baseball game has only 9 innings, but we still play to win. Pray for me.

\hfill

\texttt{Bobby Shiflett on 2014-02-16 at 14:26 said:}

francismercuri on 2013-08-29 at 15:09 said:

``Rabelais' name has also been associated with a short tract on Spirit evocation:

\url{http://www.sacred-texts.com/grim/moses7/m7112.htm}''


How is the Tabellae Rabellinae associated with Rabelais? By whom? The (highly doubtful) claim is that the Rabellini was published in 1501. (Though some scholars place Rabelais' birth to have been in 1483, the larger consensus is that he was actually born in 1494) In one scenario he would have been 18, the other, 7. There is no association.

\hfill

\texttt{francismercuri on 2014-02-17 at 16:17 said:}

Bobby Shiflett, I can't recall ``who'' had the theory, which is why no reference was included--in retrospect, agree though, spurious at best, or no relationship at all.

\hfill

\texttt{Cologero on 2017-12-14 at 09:52 said:}

This is Evola's essay on Crowley from his book on contemporary spirituality:

\textit{Evola on Crowley}\footnote{\url{http://www.gornahoor.net/library/EvolaOnCrowley.pdf}}

\hfill

\texttt{Caveman on 2017-12-14 at 11:54 said:}

Evola in his youth consorted with satanists like Maria de Naglowska; it seems that such ventures left a permanent mark on his soul, so he ever remained biased against the Christic Revelation and many of his perspectives were distorted by subjective errors. This is again demonstrated in his absurd approval of Crowley. That this site can endorse his assessment of Crowley, who was an open enemy of Christianity and whose `revealed' Book of the Law preached the destruction of all real Tradition, and spat on both Jesus, Mohammad and Buddha (and called compassion a `vice' and exalted the power-ethos of satanism) is incredible. If Crowley and his `revelation' were not instruments of counter-initiation and counter-tradition (as Guénon himself was convinced), what is? The being from whom Crowley channelled the book and became a disciple was not some harmless ``sprite'', but a demonic servant of the Enemy, and his satanism was not a mere ``facade''.

To take the destruction of his associates this lightly is disturbing to say the least. If nothing more it should at least prove that he was not qualified to serve as master for others, and he will have to Pay for this presumptuous pretension and charlatanry. In reality Crowley felt no remorse for leading people into the mire; he even bragged about having destroyed people. How will your students be equipped to deal with the subtle counter-tradition to come if you have taught them to ignore the marks of the devil even when they are so explicit as in the case of Crowley? This reckless `discernment' is dangerous to say the least.

Please excuse me for the indignant tone. I usually have the highest respect for your work, but there is enough whitewashing and normalisation of Crowley as it is elsewhere. Sites such as this one should help people discern why Thelema does not represent a valid esoteric path, not contribute to the chaos that dominates in the European occult scene or give even the smallest recognition of legitimacy to these subversive parodies of spirituality.

\hfill

\texttt{Caveman on 2017-12-15 at 12:43 said:}

I see now that I over-reacted and misunderstood your intention. I responded a little too quickly.

There are valid insights to be found in Crowley for sure, but they are mixed into a framework that was psychically toxic and chaotic. What Evola says of the chaotic psychic source of the Book of the Law is not confined to that one `revelation' (entity-channelling, apparently accepted by Evola as genuine)---this Book is proof of the fundamental nature of Crowley's occult career, as he himself considered it absolutely essential to his mission, and believed that the Book would eventually annihilate the Christian civilization (and other traditional orders worldwide) and usher in the New Aeon. He even believed that the power behind that Book constituted the occult cause of the tragedy of the two world wars. I postulate that it is the valid esoteric elements in Crowley that represent the actual facade, while it is the satanic spirit that was the driving power at the centre of Crowley's magickal career. You and Evola got this the wrong way around; Guénon, on the other hand, was as usual more clear-sighted. It is true that one shall not judge by superficial outward appearances, that satan often transforms into the appearance of an angel of light, or a wolf in sheep's clothing. (It is not impossible that even a few of the more corrupt and evil Popes were consciously subversive agents, for example the case of de Medicis effectively buying the seat of Peter.) But was Crowley really a sheep in wolf's clothing? In Christian Europe real satanists (one example being Francis Bacon: I refer to your article regarding him and the creation of modernity) have in general concealed their true identity, they themselves operating in a manner analogous to the hermetic tricksters that you mention. But times have been changing, meaning that such forces have been able to reveal themselves more openly, being necessary in order to normalise certain ideas and phenomena. Pure satanism, the mysteries of darkness and of the fallen angelic hierarchies, cannot manifest itself exoterically or in public as these would make little more sense than Chaos to those whose human consciousness has not been forced open to the Abyss and been effectively `possessed' by the influences in question. These mysteries cannot be revealed in their pure form publicly because they are completely anti-human. The various popularized exoteric versions of satanism can only serve as preparation for real satanic initiation, which necessitates gradual demonic possession of the conscious being through destruction (relatively speaking of course) of the human heart and soul. So in Crowley the satanic spirit was mixed and dilluted with elements of real esoteric spirituality for human beings, but it was still present and indwelling for all that. If the marks of the devil he flaunted quite openly even on his `facade' were meant to repel certain types of men, who would not be repelled? Not only those who are attracted to satanic trappings, but also those sincere seekers who do NOT take satanism seriously as a reality to begin with, who think it is just some superstitious Abrahamic invention, and who are therefore at risk of being exposed to spiritual evils they are not prepared to fight because of their ignorance, which they are deluded to see as intelligent and sophisticated. No authentic esoteric initiate within the Christian European civilization would take the satanic as lightly as Crowley did, simultaneously denying out of hand any reality whatsoever to the Christian conception of the Devil, the genius of evil, and praising the ancient Shaytan as understood by himself. Here there is either cunning manipulation at work by one who knew more of real satanism than he would us believe, or again simple ignorance by one who cannot thus be a real initiate but a charlatan whose only domain of real direct knowledge (that is not just abstract philosophising) is that of sorcery, lacking even full initiated into the `counter-initiatory' mysteries, or those of `black magick'.

To acknowledge Crowley as a sorcerer of a certain effective power and a man of some theoretical knowledge that is valid and applicable would be in accordance with the truth of the matter, but it does not follow from this that he was also a man of true spirituality, or one qualified to be a master leading others on the straight and narrow path to God. When dealing with a phenomenon like Crowley one cannot, when discerning the spiritual nature and function of the man, choose to focus on parts of his writings that are valid in isolation as if these concern a different Crowley than that represented by the totality, by the unifying principle guiding his agenda. If a certain idea reflects the Good and the True, but the man's activity as a whole goes counter to the Good and the True, then this idea has only an illusory relationship to the substance of this men's activity in the world.

By their fruits ye shall know them. Which side of the spiritual war did Crowley's work actually benefit, whether through his ideas, his living influence on his associates, and through the work carried on by his direct personal disciples? The ultimately disintegrating and chaotic influences working through Crowley should be evident and can be seen directly, without having to navigate a web of complex ideas to find the ``real'' identity of Crowley hidden by a contradictory facade. In Crowley's case, the driving force or directing intelligence of his mission corresponds truly with the satanic aspects of his image, and it is the purer spiritual ideas found scattered throughout his writings that become maya, illusion obscuring his essence. To view Crowley the man through the lens of legitimate spiritual ideas is to let down one's guard and possibly become subtle influenced through this sympathetic opening by the dark forces attached to the man and his work. The hermetic way is to see such men directly, intuitively, in terms of depth of forces at work and breadth of influences emanating from them in the world. Evola understood certain things about the reality of Crowley's work, such as the origin of the Book of the Law, but failed to see how the same phenomenon applied to Crowley himself, who was the individual

`incarnation' of the same influence manifesting itself in that Book, and it showed in the most critical aspects of his life, actions and relations; this is the fundamental reality of the man and his mission, and the valid esoteric elements found scattered throughout his body of writings become more or less accidental or of shallow significance in view of that. Crowley appears so `complex' because he was chaotic, but this complexity does not necessarily conflict with a simple and singular essence in terms of a unifying infernal will expressing itself through that chaos. What can be said of truth is that it essentially manifests itself through initiates in simplicity, order, beauty and purity, even though it may take on unexpected and unconventional appearances. What it never does, even in antinomian expressions, is chaoticize people who encounter it, and it hardly wants children to take part in orgiastic sex rites.

Did really Crowley die with ``all his faculties intact''? What were all these facilities that were so intact? He was a drug addict, but some have said he needed it for medical reasons. In terms of a superficial modern understanding of man, perhaps he had ``all his faculties'' intact, but ther le are even ordinary human `faculties', such as empathy or conscience, that can be lost without this being necessarily outwardly visible as impairing `functionality', though it may often be inferred from their manipulative and egotistical way of relating to others. Even profane science has been able to empirically prove that there are psychopaths whose hearts are basically dead, even though, due to their often times great skilfulness in navigating society and attaining power and influence through great `intelligence' and cunning, they may be able to give a superficial appearance of having ``all their faculties'' intact. Hillary Clinton, for example, is certainly such a case, but they are very, very common, not least in the higher echelons of society. Psychiatrist/psychologist M. Scott Peck in his book `The People of the Lie: Hope for Curing Human Evil' expressed his view that individuals in whom the typical dramatic symptoms of demonic possession are visible are actually resisting a partial evil possession, while those who are perfectly possessed without resistance will not be visibly dysfunctional in such a way. Fr Malachi Martin said something similar. The people who collapsed around Crowley may have been more human than himself, while his own functionality until the end might have been becaus he had managed to separate his demonised ego from his human soul, having partly sustained himself by `vampiric' exploitation of others. These things are real, and not taking them into account when analysing cases like Crowley leads to limited insights. It is important to understand that the particular faculties that appeared to be intact in Crowley prove nothing insofar as the question of his salvation or damnation is concerned, nor does his display of a certain conceptual `knowledge', and even less his magickal powers. A humble life expressed through love (a word that for Crowley didn't mean much more than f***ing in practice) would have been a much stronger proof than any of his grandiose magus ego trip. That was not merely a hermetic mask, because the real hermetist chooses a humble mask, one that does not represent a vain and corrupting example for others. It is not in my place to judge anyone for the paths that they have chosen, but when someone makes himself a public influence and attempts to revolutionize the world as a false prophet (which he most definitely DID intend), we have a problem.

\hfill

\texttt{Andrea (Oxford) on 2017-12-15 at 20:37 said:}

Crowley was a good magician, but a charlatan from an authentic initiatic perspective. Once you study in depth the symbolic structures of the Tree of Life that Crowley already used in A?A? for the progression of the degrees, than you realise how he abused the ``sphere of Yesod'' because of his altered states of consciousness due to drugs consumption and sexual perversions which have been summarised by Gianfranco De Turris and Vittorio Fincati. If you want to master magic art, the knowledge of Crowley's work is relevant but let's not call it Initiation.

I am very reluctant to talk about magick in general as I am fully immersed in the ordering principles of the Way and I am no longer interested in magical techniques, having identified in the Alchemic Art the only corpus of procedures coherent with metaphysics and the traditional approach.

Marco Barsotti has recently came out with a well documented book on the mythical story of the Babylonian Goddess Ištar which seems to be quite interesting... \hfill

\end{sffamily}
\end{footnotesize}

